% Generated draft: requires mathematical and source review.
\documentclass[11pt]{article}
\usepackage[margin=25mm]{geometry}
\usepackage{fontspec}
\setmainfont{Latin Modern Roman}
\usepackage{amsmath,amssymb,mathtools,unicode-math}
\setmathfont{Latin Modern Math}
\usepackage{enumitem,needspace,xurl,hyperref}
\setlength{\parindent}{0pt}
\setlength{\parskip}{5pt}
\setlength{\emergencystretch}{3em}
\newcommand{\N}{\mathbb{N}}
\newcommand{\Z}{\mathbb{Z}}
\newcommand{\Q}{\mathbb{Q}}
\newcommand{\R}{\mathbb{R}}
\newcommand{\C}{\mathbb{C}}
\newcommand{\F}{\mathbb{F}}
\newcommand{\A}{\mathbb{A}}
\newcommand{\PP}{\mathbb{P}}
\newcommand{\E}{\mathbb{E}}
\newcommand{\eps}{\varepsilon}
\newcommand{\dd}{\,\mathrm{d}}
\begin{document}
\section*{UnsolvedMath ID: 700765}
\textbf{Source problem code:} catalogue entry
\subsection{Definitions and mathematical statement}
A variety of groups \(\mathcal V\) is the class of all groups satisfying some specified collection of group identities. A group \(G\in\mathcal V\) is relatively free in \(\mathcal V\) if there are a set \(X\) and a map \(\iota:X\to G\) whose image generates \(G\), such that for every \(H\in\mathcal V\) and every map \(f:X\to H\), there is a unique homomorphism \(\varphi:G\to H\) with \(\varphi\circ\iota=f\).

A group is finitely presented if, as an ordinary group, it admits a presentation \(\langle S\mid R\rangle\) with both \(S\) and \(R\) finite. A group \(N\) is nilpotent if its lower central series, defined by \(\gamma_1(N)=N\) and \(\gamma_{j+1}(N)=[\gamma_j(N),N]\), satisfies \(\gamma_{c+1}(N)=\{1\}\) for some integer \(c\geq 0\). Here \([A,B]\) denotes the subgroup generated by the commutators \(a^{-1}b^{-1}ab\), with \(a\in A\) and \(b\in B\). A group is virtually nilpotent if it has a nilpotent subgroup of finite index.

\textbf{Question (A. Yu. Olshanskii, 1990).} For every variety of groups \(\mathcal V\) and every group \(G\) relatively free in \(\mathcal V\), does finite presentability of \(G\) imply that there exist a subgroup \(N\leq G\) and an integer \(c\geq 0\) such that
\[
[G:N]<\infty
\qquad\text{and}\qquad
\gamma_{c+1}(N)=\{1\}?
\]
The terms ``relatively free'' and ``finitely presented'' are interpreted here in their standard senses defined above; the supplied record gives no additional restrictions on the variety or the free generating set.
\subsection{Short English statement}
Must every finitely presented relatively free group contain a nilpotent subgroup of finite index? The question is attributed to A. Yu. Olshanskii (1990). The supplied registry labels it open; that status has not been independently verified.
\subsection{Sources}
\begin{itemize}
\item Local UnsolvedMath registry, permanent ID \(700765\), source problem code \texttt{catalogue entry}, title \emph{Kourovka Notebook Problem 11.73}; supplied upstream revision \texttt{d371ee917ac9647fbb06cd6fc529deec6f3dda82}.
\item \emph{The Kourovka Notebook (Unsolved Problems in Group Theory)}, 11th issue (1990), Problem 11.73; proposed by A. Yu. Olshanskii. The supplied background also identifies the 21st issue edition (2026), \texttt{arXiv:1401.0300v48}.
\item The supplied collection description names E. I. Khukhro and V. D. Mazurov as editors of the Kourovka Notebook collection covering Issues 1--21 (1965--2026). The bibliographic information and open-status label are supplied source data, not an independently verified literature assessment.
\end{itemize}
\end{document}
